% Derived from the owned full Markdown narrative; source-stage only.


\section{Assistance and Retained Learning}

Let a computer guide a hand through a tune, then stop the teaching commands.
What can the trainee play now? The first performance shows what the combined
system could do; the second begins to test what the nervous system retained.
Any movement bridge needed to express the trainee's own control must remain
available in both comparisons. Otherwise we would be testing a disconnected
hand and calling the result a test of memory. Brain2Body connects
nervous-system measurements to physical outcomes to make this distinction
testable. The numeral 2 means ``to''; body-to-brain feedback matters as much
as the outward command.\par

Rapid acquisition of difficult physical skills is the ambition. Sewing, playing an instrument, manipulating a puzzle and coordinating a whole body put that ambition under pressure: each demands more than one correct movement. The programme must explain how useful guidance could become a capability the trainee retains. These are proposed tests, not demonstrated abilities. Neither arbitrary skill upload nor complete nervous-system reproduction is established, and the current models have not passed their original performance gates.\par

A decoder turns an observation into a prediction. Before calling it correct, we must decide what answer it is meant to give. NeuroTranslation proposes relating neural signals and their representations between explicitly described sources and targets; Brain2Body follows that prediction into a body, where movement changes the feedback available for the next action and learning becomes a further question. The argument therefore needs an independently measured action, a model whose contribution can be isolated, and a retention test that can fail. Published research motivates those requirements. The companion computational work supplies bounded comparisons, including negative results; it has not supplied evidence of biological teaching.\par

\section{Compare Predicted and Measured Movement}

Perception is subjective. Movement is not. A person experiences a sensation; an independent instrument can measure a finger's position, a contact, a force or the timing of a release. That contrast does not make every movement measurement exact or reveal private experience. It gives the programme a tractable starting point: relate an uncertain neural observation to an outcome that can be checked.\par

Consider a key press. Choosing the right key, reaching it, pressing at the right time, releasing it and avoiding neighbouring keys are separate outcomes. A video can show position; a contact sensor can show a strike; a force sensor can measure force. A task score records success under declared rules. Neural activity near an electrode is another observation. None silently supplies the others.\par

In the proposed loop, sensory input guides action. The body and environment
change; measured outcomes supply feedback for the next action. Feedback control
adjusts action using the difference between the desired and observed outcome.
Todorov and Jordan provide an established account of motor coordination under
uncertainty. Applying it here is a programme rationale, not an empirical result
of their paper. Every claimed loop must disclose its delays and whether its
feedback came from a living subject, a simulation or a recorded replay.
\href{https://www.nature.com/articles/nn963}{Todorov and Jordan, 2002}.\par

Accuracy means agreement with an independent reference. Precision means repeatability under stated conditions. A consistently wrong model can be precise. Ground truth means an independently acquired measured reference with its own calibration, sampling limits and uncertainty; it is not error-free access to biology. Finer timesteps, more decimal places and a confident animation cannot establish these properties. \href{https://www.nist.gov/pml/nist-technical-note-1297/nist-tn-1297-appendix-d1-terminology}{NIST measurement terminology}.\par

\section{Structure, Activity and Function}

A movement score tells us what happened. Explaining why requires more than a map of the routes a signal could take. Timing, the cells' responses to chemical signals, the body's current state, learning history and task conditions affect what follows. An intervention deliberately changes something to test its effect; a connection map alone cannot tell us that response. Structure narrows the explanation, while measured function must test it. That is why a body-grounded prediction needs both.\par

The Human Genome Project offers a bounded analogy. Obtaining a reference sequence did not finish the investigation of regulation and function; subsequent projects investigated functional elements. That is a distinction between a map and its interpretation, not a claim that genome sequencing failed. \href{https://www.genome.gov/11009066/2003-release-scientists-venture-deeper-into-the-human-genome}{NHGRI, Beyond Genes}.\par

The MICrONS Consortium combined recorded visual responses with detailed connectivity in a volume of mouse visual cortex. Its strength is the combination of different measurements. It does not provide a whole human nervous system, peripheral or spinal signals, or an unrestricted account of learned skills. Co-registration means relating measurements to corresponding locations; such correspondence must be demonstrated for the actual source, not borrowed from an attractive template. \href{https://www.nature.com/articles/s41586-025-08790-w}{MICrONS Consortium, 2025}.\par

Measurement itself needs explanation. Functional magnetic resonance imaging, or fMRI, measures an indirect blood-related signal over time. In the experiments of Logothetis and colleagues, that signal related more closely to local field potentials, aggregate electrical activity near a recording site, than to spiking output. It is not a photograph of thought. Tang and colleagues demonstrated constrained semantic reconstruction with substantial dependence on the subject and cooperation; this is not universal one-to-one decoding of subjective experience. \href{https://pubmed.ncbi.nlm.nih.gov/11449264/}{Logothetis et al., 2001}, \href{https://www.nature.com/articles/s41593-023-01304-9}{Tang et al., 2023}.\par

Quantum measurement can illustrate that observing a system may involve interaction. It does not establish quantum backaction, the influence of a quantum measurement on a system, as the cause of these large-scale neural limits. Scanner constraints, indirect signals, person-specific learning and ambiguous representations supply concrete explanations that must be tested on their own terms. The analogy stops before it becomes an unsupported brain mechanism. \href{https://www.nobelprize.org/uploads/2018/06/advanced-physicsprize2012-1.pdf}{Royal Swedish Academy, quantum measurement background}.\par

\section{Existing Brain--Body Interfaces}

Bouton and colleagues linked recordings from motor cortex, movement-related tissue in the brain's outer layer, to electrical activation of forearm muscles, enabling useful hand and wrist movements in one participant with quadriplegia, paralysis affecting four limbs. This demonstrates a participant-specific route from brain signals to bodily action. It does not demonstrate an unfamiliar skill transplant or a complete map of the nervous system. \href{https://www.nature.com/articles/nature17435}{Bouton et al., 2016}.\par

Lorach and colleagues linked cortical recording to spinal stimulation and evaluated standing and walking in a participant with chronic spinal injury. They also reported rehabilitation-associated improvement with the interface switched off. Movement in a real participant provides functional evidence beyond a rendered demonstration. The study remains a particular participant, intervention and rehabilitation history; assistance, retained improvement and tissue regeneration are distinct claims. \href{https://www.nature.com/articles/s41586-023-06094-5}{Lorach et al., 2023}.\par

Flesher and colleagues used stimulation of sensory brain tissue to provide touch-related feedback during robotic arm control. The work supports testing the return path of an interface. It does not establish transfer of an expert's repertoire. These studies motivate bounded experiments; they cannot be stacked into proof of a new integrated system whose relevant links have never been measured together. \href{https://pubmed.ncbi.nlm.nih.gov/34016775/}{Flesher et al., 2021}.\par

\section{Digital-Fly Models}

Digital flies sharpen the question of attribution: when a composite system
moves well, which component supplied the ability? The answer determines
whether the result tests neural wiring, a trained controller, body mechanics
or their engineered combination. Those contributions must remain visible
before a movement is credited to a reproduced nervous system.\par

FlyWire reconstructed 139,255 neurons, nerve cells, in one adult female fly brain. This is an exceptional structural resource with reconstruction and coverage limits. A map of chemical connections does not supply every connection's physiological effect, the complete nervous system, the body or the environment. Pospisil and colleagues propose combining intervention measurements with structural constraints to estimate causal effects; their estimator tests in simulation are not a completed empirical map of all fly-brain dynamics. \href{https://www.nature.com/articles/s41586-024-07558-y}{Dorkenwald et al., 2024}, \href{https://www.nature.com/articles/s41586-024-07982-0}{Pospisil et al., 2024}.\par

Shiu and colleagues show the positive route: specific connectome-derived predictions were tested through neural stimulation and behaviour. A connectome is a map of neural connections; it supplied constraints, not an automatic functional certificate. Useful predictions coexist with missed experimentally observed Phantom and Usnea effects and incorrect G2N-1/Usnea water-silencing predictions. The authors report 91\% agreement across 164 tested predictions, or 84\% after excluding one mostly negative screen. Neither figure measures whole-fly fidelity. Missing receptor dynamics, electrical connections and chemical modulation limit the account, including absolute simulated firing rates. \href{https://pmc.ncbi.nlm.nih.gov/articles/PMC11446845/}{Shiu et al., 2024}.\par

Body mechanics introduce another layer. NeuroMechFly v2 documents simulated leg tips slightly below the ground because of numerical errors or physics instability, with a 0.05 mm margin to prevent false overstretch detection. Its adhesion model uses engineered forces and stepping phases while biological liftoff mechanisms remained uncertain. These are specific disclosed modeling limitations, not evidence of dramatic tunnelling or an ability no living fly could ever acquire. FlyBody uses trained artificial controllers to imitate measured movement; its time-averaged aerodynamic approximation and omitted detailed muscle and sensory physiology remain separate limits. \href{https://gizemozd.github.io/assets/pdf/2024_neuromechflyv2.pdf}{Wang-Chen et al., 2024}, \href{https://pmc.ncbi.nlm.nih.gov/articles/PMC12310536/}{Vaxenburg et al., 2025}.\par

Eon's March 2026 technical account describes hand-chosen brain--body mappings, substantial reliance on existing body controllers, largely missing state and learning mechanisms, and internal dynamics not yet validated against relevant biological signatures. That dated developer disclosure supports a conclusion about an engineered integration, not a complete biological upload. We have not independently reproduced its dynamics. \href{https://eon.systems/updates/embodied-brain-emulation}{Eon Systems, 10 March 2026}.\par

The same attribution problem appears in games. Pinned fly-chess code supplies symbolic board information, selected mathematical dynamics, a trained readout and external legal-move search. Pinned DoomFly code uses the distinct male central-nervous-system map and injects game-engine enemy geometry into model neurons after the intended visual route reportedly failed. These observations come from inspected code and author accounts, not reproduced gameplay or evidence that a living fly learned either game. \href{https://github.com/mncrftfrcnm/making-fly-play-chess/blob/b97dfefbf9b8850f3cdf6b412c2ed840be94e260/scripts/fly_chess_trainer.py}{Fly-chess trainer}, \href{https://github.com/shreyash-sharma/doomFly/blob/5f6c2881be5f5f40b0592a75e6ee231b50333b40/src/doom/fearless/danger.py}{DoomFly threat input}.\par

\section{Simulation Failures}

Attribution also depends on the world in which the controller succeeds. A
controller can improve against a score while moving farther from the intended
task. To use simulated practice as evidence, we need to know both what the
score rewards and whether the simulated movement obeys the relevant
physical constraints.\par

An optimizer, the procedure adjusting a model to improve its objective, can exploit an incorrect simulator or a poorly chosen score. First-hand historical accounts describe virtual creatures gaining propulsion from numerical errors and soft robots receiving artificial energy from ground-penetration corrections. CoastRunners illustrates a different failure: circling reward targets improved the score while leaving the race unfinished. Physics violation and objective misspecification need different diagnoses. Neither example establishes that every fly model shares those failures. \href{https://arxiv.org/html/1803.03453v4}{Lehman et al., first-hand accounts}, \href{https://deepmind.google/blog/specification-gaming-the-flip-side-of-ai-ingenuity/}{Google DeepMind, specification gaming}.\par

Protein prediction offers another carefully bounded comparison. AlphaFold3's original paper describes invented structure in unstructured regions and overlapping predicted chains. Its antibody-interface benchmark reports 62.9\% cluster-averaged success at DockQ > 0.23, a score threshold for agreement with experimentally determined interfaces, hence 37.1\% failing that threshold. It covers 65 interface clusters from 71 complexes, with confidence-based selection from 1,000 seeds and five samples per seed for each target. This is not a universal percentage of proteins that fail to function. Crystallization is one measurement bottleneck; other methods include cryo-electron microscopy, which images samples using electrons, and nuclear magnetic resonance, which uses atomic responses to magnetic conditions. Synthetic examples can help training, but they cannot create independent biological validation of unmeasured function. \href{https://www.nature.com/articles/s41586-024-07487-w}{Abramson et al., 2024}, \href{https://www.nature.com/articles/s41586-021-03819-2}{Jumper et al., 2021}, \href{https://pdb101.rcsb.org/learn/guide-to-understanding-pdb-data/methods-for-determining-structure}{RCSB structure methods}.\par

Biological virtualization means reproducing specified biological behaviour and causal dynamics in a digital model within independently tested conditions and error bounds. Invented biology presented as faithful reproduction is hallucination. A disclosed approximation can still be a useful scientific model. Missing validation does not prove a prediction false; one correct guess does not validate its mechanism or general reliability. The dice analogy concerns unknown reliability, not equal odds or an assertion that learned predictions are random.\par

\section{Recording Access, Context and Stable Signals}

The proposed measurement chain has to survive three questions. Access it: can we obtain the nervous-system observations relevant to a defined question, and what burden does access impose? Ground it: can we relate those observations to independently measured body behaviour on known clocks, in stated quantities and with stated uncertainty? Keep it: does the recording and its mapping remain useful as tissue, movement and signal quality change? Keeping the recording usable is a different endpoint from a trainee retaining a skill.\par

Sources need organism, subject, region, task, modality, units, sampling and acquisition conditions. Cortex, spinal cord, peripheral nerve, muscle and movement are different observables. Electroencephalography (EEG) measures electrical activity using scalp electrodes; electrospinography (ESG) records signals over the spine; electroneurography (ENG) records nerve electrical activity; electromyography (EMG) records muscle electrical activity. EMG is not direct peripheral nerve recording. Separate studies of these signals do not become one synchronized voluntary cortex--spine--periphery chain by placing their charts together.\par

The resources associated with HHMI/Janelia, Allen, Harvard and other groups contribute particular observations to particular questions. Their value depends on that match. A broad catalogue cannot supply a paired signal that was never recorded; a model can estimate the missing value, but the estimate remains a prediction. An independent measurement can check it, but must remain a separately identified observation. When data are combined across organisms, source identities and physiological differences must remain traceable. Whole-body coverage is the objective to investigate, not a property conferred by the catalogue.\par

\section{Physical Constraints}

MyoSuite supplies muscle-based simulations with contact and task objects, making it possible to examine how a command becomes movement. The proposed chain keeps the neural predictor, goal, muscle controller, body mechanics and outcome score distinct. If the predictor selects the wrong key, a competent controller can execute that mistaken goal perfectly. Setting the joint directly to the desired pose would bypass the muscle-control question and must not be reported as a learned muscle action. These controlled computational studies still need their own tests of participant-specific biological accuracy. \href{https://proceedings.mlr.press/v168/caggiano22a.html}{Caggiano et al., 2022}.\par

NeuroTranslator names the implemented software joining particular recordings, models, replay and computer-body checks. NeuroTranslation Network names a proposed design for joining further components while keeping their origins and meanings clear. Neither name establishes full nervous-system coverage. Model size, learning from neural recordings and learning through simulated rewards require matched comparisons on the actual task and available hardware. A larger model or an impressive reward does not replace missing measurements.\par

Reaching, contact and release offer possible task relationships to compare across organisms. The mechanics still belong to each organism: a hand, paw and whole-body quadruped have different proportions, joints and control requirements. An averaged anatomy cannot establish a mapping that works for any of them. Physical scale, computational detail and statistical aggregation must also remain separate; a small organism can have a detailed model, while a pooled error can hide a failed recording. Biological generalization requires a test in the organism and conditions named by the claim.\par

\section{Current Study: Neither Model Passed}

The current hand study puts that demand for comparison into practice. Does
retaining recent neural history improve the stated prediction without
violating the other required conditions? Finishing a training run cannot
answer it; the original conjunction of tests must decide.\par

The current development study comprises six sessions from one macaque: four training dates and two development-validation dates. Ninety-six recording contacts in the source's HandM1 motor-cortex region contribute threshold-crossing-rate (TCR) and spiking-band-power (SBP) measurements, counts of detected signal crossings and power in an archived signal band. They give 192 features, not 192 independently located electrodes. Its targets are two source-normalized finger-sensor groups: index and combined middle/ring/small fingers. They are not measured joint angles, forces or a full-body state.\par

Eighteen runs saved 720 epochs, passes through the training material. Both gated recurrent unit (GRU) candidates, sequence models with a mechanism for retaining earlier information, fail their original model gates. Both fail the required no-harm and positive neural-information conditions: the predeclared comparisons must show the required absence of harm and contribution of the neural information. All five separate history, order-use and auxiliary-benefit claims also FAIL. Reduced aggregate error relative to a selected ridge baseline, a simpler predictor trained with a penalty on large fitted coefficients, does not override a failed required combination of criteria. The negative report records these failures; it does not establish that either model is effective.\par

The failed comparison also leaves a reproducibility problem. Both candidates use Torch automatic differentiation to calculate derivatives and Adam to update the model from those derivatives and its stored state. Agreement between candidates using that common route cannot independently certify the route. Full independent NumPy calculation using a distinct numerical software library remains unverified, as do complete model/update state, random-number-generator (RNG) state, reload and training-trajectory evidence. The final-checkpoint/update gate therefore remains unresolved.\par

A continuation would need its own verified execution clearance, which is also currently unverified. The four reserved dates remain CLOSED. This article reports no new fitting, optimizer execution or reserved-test evaluation, and neither an accepted model nor full-training graphics-processor (GPU) promotion. The saved failures remain the result to explain. \href{research-evidence.json}{Programme computational evidence summary}.\par

\section{Earlier Studies}

An earlier human finger-attempt classifier obtained 116 correct predictions among 150 later trials after calibration using the first 50 labelled examples of each recording session. In its separate keyboard replay, 41 of 60 source-labelled right-finger trials produced the intended contact, while all 54 supported commands actually issued were executed; 13 presses were unintended or mistimed. The controller's execution score does not fix a prediction error. The model waited for a completed window; its earliest decision was one second after the cue plus computation. Inspected trials become development material if later tuning uses them.\par

A separate goal-driven simulated policy completed 32 of 32 scored phrase cases: eight reserved simulated orders crossed with two tempos and two poses, not 32 independent organisms. It uses a constrained generic four-key hand whose wrist, thumb and sideways finger-motion joints remain braced. It receives requested keys, phase, duration and smooth fingertip goals; displayed neural clips do not drive the simulated hand. The record preserves an interrupted original evaluator and a completed independent replay check; no test was reopened. Soft simulated key limits are not guaranteed physical travel bounds. A current saved development replay contains 336 frames at a 20 ms control interval, with 22 of 22 correct strikes, six of six chords, 22 of 22 timely releases, 34 ms worst timing error and no extra strikes. These are bounded simulation/readout results, not biological teaching, neural policy learning or unexpected real-world musical mastery.\par

The human recordings, macaque hand mapping and separate human EEG/ESG/ENG/EMG clips have distinct acquisition conditions and clocks. They do not demonstrate a voluntary central--spinal--peripheral pathway or drive a combined macaque/human arm simulation. Actual joint and force ground truth and physical unit calibration remain gaps. A task label, a normalized sensor value and a simulated joint angle must retain their different meanings. \href{research-evidence.json}{Programme computational evidence summary}.\par

\section{Failure Analysis}

The failure narrows the next question: where did the extra error come from? In the reported index-finger comparison, the auxiliary decoder's excess mean squared error, the average squared prediction difference, separates mainly into a mean offset and a smaller difference in residual variation. That decomposition locates a mismatch without identifying its cause. Trial composition, recording changes and state-dependent offsets could each contribute. Re-centering the validation predictions after inspecting that miss would use the answers to repair the prediction, so the adjusted result would no longer preserve the independent test.\par

Training-only context summaries show variation within trials, between trials and between dates; correlated overlapping frames are not independent observations. Named scalar calculations have bounded independent checks, including a preserved rejected skewness calculation, a measure of distribution asymmetry, and corrected arithmetic. Those checks do not demonstrate physiological state, neural drift or a useful new decoder. The proposed small context branch remains unexecuted. Any next fit requires independently checked training-only selection, normalization, controls, initialization, uncertainty and full update/state evidence, while preserving every original threshold. \href{research-evidence.json}{Current evidence and limits}.\par

\section{Viewer Capabilities and Missing Measurements}

A graphical user interface (GUI) should let the reader tell a measurement from a prediction, a simulation from a missing value. Saved replay must preserve the source-time clock and complete-run metrics, even when the reader pauses at a persuasive frame. Cache-only playback displays an identified saved result without silently starting a new experiment; an uncached selection must show absence, and a failed source must clear stale geometry and results. Natural-language goals should become explicit task specifications. Receiving files in a chat does not itself establish inference or body control, so the interface must show what operation the system actually performs.\par

A saved replay and a live research view therefore need separate checks. The reader should be able to see which one is available, which source it uses and which measurements are missing.\par

\subsection{What the online and offline versions require}

The online candidate needs its whole directory. A disconnected reader instead
needs the separate static export compiled from the original source; research
views that use the application programming interface (API), the interface to
the software service, are unavailable offline. Neither route has current
whole-interface verification. Earlier bounded replay checks cannot establish
that the present interface or the separate export works.\par

\paragraph{Reproducibility record.}
Moving opaque resources from the single HTML web-page source file into the
package reduced that file from 113,298,857 to 34,295,597 bytes; its assets
remain required for whole-directory online delivery. The separately compiled
original-source offline export remains a different artifact.
Browser verification was unavailable, and the candidate remains unverified and
unpromoted. These are packaging measurements: they establish no interaction or
loading latency, memory, crash-resistance or GPU benefit and no newer browser
qualification. \href{research-evidence.json}{Computational evidence and limits}.\par

\section{Anatomical and Functional Validation}

A viewer can expose an error only if the body it shows is itself accountable to references. Begin with proportions, muscles, bones and joints, then trace the neural routes and check contact and viewing geometry wherever those features are claimed. Each component and the assembled scene need lifelike three-dimensional review against a series of real-world references. Biological identification colours, useful close-ups and tilted orthogonal views can make those comparisons legible. Transparent layers, cross-sections and paired views can reveal relationships; illustrative signal paths must remain distinct from physical parts. This review tests the anatomy we show. Hidden physiological function needs its own evidence, however plausible the image. A conceptual schematic remains an explanation, not a biological acceptance result.\par

A template surface or atlas provides a reference. Registration has to establish its correspondence to the recorded individual's anatomy before a location in the template can be attributed to that individual. The named source-binary reconstruction matches template geometry and atlas values, and its optional region boundary retains disclosed surface-mesh defects. Those bounded correspondences neither localize the recorded neurons nor establish their function. The retained framework illustration is historical and conceptual; its labels describe earlier computational work, not the current model-gate outcome.\par

The proposed left and right tripod views of the hands and arms need calibration, occlusion checks and independently verified timing. The motion-capture relay-composite would align or stitch those views. A separate perpendicular LiDAR view, using light to measure distance, may make particular surfaces easier to observe; it is not guaranteed to reveal every surface or solve neural registration. The acquisition method must earn its reconstruction claim against real measurements.\par

\section{Proposed Skill Tests}

Once the movement and its measurement can be trusted, the next task should expose something the present test cannot answer. Key presses separate selection, contact and timing errors. Threading adds flexible-object geometry, friction and corrective contact; puzzle manipulation adds object state and coordinated grasping. Archery or breakdancing needs forces, balance and wider body coverage absent from the current hand study. This gives the task ladder a reason to grow without pretending that a classifier or successful replay has already climbed it.\par

Reserve new task orders, object properties, starting positions and recordings before selecting the method. Then give simpler controllers and copied-example methods the same information and practice, and score the complete task: success, unintended actions, delay and failure. A selected clip cannot answer that comparison. Right-angle and non-right-angle mazes could test orientation and memory demands, but their difficulty needs measurement; 10\ensuremath{\times} or 100\ensuremath{\times} labels do not provide it. Computer training duration likewise cannot establish biological learning duration.\par

\section{Development, Twins and Disease}

The body also changes while a mapping is being learned. The native-development branch asks how recording access, maturation and changing function affect that mapping. Cook and colleagues provide species-specific skeletal staging in the fat-tailed dunnart. It does not establish corresponding timing in antechinus, a safe neonatal neural interface or an artificial rearing system. Armstrong and Ferguson studied age-related skin-wound scarring in \emph{Monodelphis domestica}; skin healing does not certify scar-free neural access, spinal repair or stable electrodes. \href{https://www.nature.com/articles/s42003-021-02506-2}{Cook et al., 2021}, \href{https://pubmed.ncbi.nlm.nih.gov/7750642/}{Armstrong and Ferguson, 1995}.\par

Identical-twin, non-identical-twin and unrelated comparisons would ask which sources of variation affect a mapping. Similar genetics leaves differences in experience, physiology and recording location to be measured. Each comparison needs matched tasks, developmental stages and acquisition conditions before a channel or skill-transfer claim can be interpreted. The twins model must undergo component and assembled-scene reference review as well; kinship and a convincing linked image supply neither functional evidence nor visual acceptance.\par

Human-cell and graft branches ask defined questions about tissue contribution and host interaction. Revah and colleagues found maturation and circuit integration of transplanted human cortical organoids, laboratory-grown tissue representing aspects of the brain's outer layer, in rats. This does not make the whole host human, transplant expertise or demonstrate injury repair. The later image-panel correction is retained in the citation record. Native development, altered development, mixed-species tissue, grafts and disease conditions are separate proposed comparisons, each requiring its own evidence. \href{https://www.nature.com/articles/s41586-022-05277-w}{Revah et al., 2022}, \href{https://www.nature.com/articles/s41586-026-10943-4}{author correction, 2026}.\par

A disease experiment must identify the impaired function and the change produced by the interface. Compensation for interrupted motor output, improved task performance, retained learning and tissue repair require different endpoints. The cause of one condition may leave a route usable that another condition disrupts, so a benefit cannot be carried across them without its own test. These branches remain proposed; the present computational record demonstrates none of them.\par

\section{Next Experiments}

The failed comparison leaves us short of a qualified predictor, before we even reach the teaching question. The measured targets are still two normalized finger-sensor groups, without calibrated joints, forces or a complete body state, and neither recurrent candidate is admitted. A constrained controller could be studied by declaring its observations, permitted actions and feedback. If it receives explicit goals, its result must remain a goal-driven result; it cannot stand in for a neural prediction or evidence of biological learning. Only with those contributions separated can a trainer comparison ask whether useful movement assistance leaves lasting capability.\par

The next comparison can start with the grasp itself. The proposed two-person movement-only study asks whether pretraining on a qualified source person's movement improves a target person's movement forecast. Compare target-person pretraining with source-person pretraining under deliberately incorrect correspondence and under the proposed correspondence, matching architecture, initializations, update budgets, target-training batches and final-checkpoint selection. The proposed shared representation contains 13 movement features and 13 missing-feature masks: wrist-relative index, thumb and object sensor positions, digit aperture and causal wrist velocity. Proper wrist frames and a characteristic sensor-distance scale fitted only on each person's training observations separate coordinate orientation and size; a common 0.1 second reference interval also makes the velocity features dimensionless. The scale is not a measured bone length. Frames, units, clocks, gaps and sensor calibration still need qualification before this representation can be used.\par

Both homologous and mechanically analogous features remain eligible for the wider programme. This particular comparison tests a declared human sensor-landmark correspondence, not cross-species neural equivalence. Motion from another recording can supply a learned movement constraint without supplying a neural observation paired to the target person's action. Scans may constrain a specified anatomical gap, but an external scan or simulated movement cannot supply that missing measured neural pairing. Keep acquisition parents, overlapping segments and derivatives in one role before windowing or fitting statistics; development observations must not influence training scales or become independent final confirmation. This comparison remains unexecuted. Any later neural adaptation must use actual paired observations and retain its distinct input sets, causal timing, controls and protected partitions.\par

A trainer experiment needs supported interfaces and new evidence. Compare experienced, inexperienced and artificial trainers with matched task information, practice, assistance budgets and relevant burden. Include sham, delayed and unrelated-feedback controls. Improvement while guidance is present does not establish retained learning.\par

Withdraw trainer information while keeping any necessary movement bridge, practice, feedback and task access constant. Then measure delayed performance and unfamiliar variations against matched controls. Record the help that remains. If a gain fails to persist, that failure counts against learning. A retained gain still cannot establish tissue repair, restored pathways or regeneration: those claims require separate biological evidence.\par

\section{Biological and Artificial Trainers}

Brain2Brain names specified information links between nervous systems. Pais-Vieira and colleagues linked recorded activity and stimulation to share task-relevant information between trained rats. Prior training and dependence on stimulation limit what the result means: it is not instant acquisition of another animal's general skill repertoire. \href{https://www.nature.com/articles/srep01319}{Pais-Vieira et al., 2013}.\par

The proposed biological-trainer experiment compares experienced and inexperienced partners with matched practice. Genuine guidance, sham guidance without the intended information, delayed guidance and matched feedback unrelated to current need test different explanations. A simple artificial controller tests whether the neural trainer adds benefit beyond a less elaborate route. Random assignment, recording of actual help and independently scored outcomes are needed to distinguish expertise, information, timing and additional exposure.\par

An artificial trainer would supply guidance from a model trained on supported observations. Its comparison begins with what the model receives and controls. Desired goals, expert trajectories, privileged labels and sensory feedback give it different resources; supplying the correct action removes the prediction problem that a neural decoder is supposed to solve. Trainers therefore need matched task information, practice and relevant burden. A useful gain while guidance remains present can support an assistance claim even if no lasting learning follows.\par

\section{Trainer-Withdrawal Tests}

Test retention by withdrawing trainer commands while holding any necessary
compensatory movement bridge constant. A bridge lets a trainee express action
despite an interrupted pathway. Remove that bridge together with teaching and
movement may become impossible regardless of learning: the test would confound
body access with skill. Keep movement access, feedback and task conditions
comparable between the baseline and later test.\par

Measure immediate performance, delayed retention and unfamiliar variations separately. Record how much help remained, whether task performance improved without trainer information, and whether the benefit exceeds matched controls. If the gain disappears when trainer commands stop, report assistance. If a gain persists, test its domain and uncertainty. Tissue repair still requires separate biological evidence; a retained task gain or successful compensatory interface cannot establish regeneration.\par

BrainNet is a network metaphor for specified links. It is not an operating market for biological mastery. Its value would have to appear in what a trainee can do through those links. Broad enablement and possible skill services motivate the programme, but they do not establish deployability, ownership of biological information or superiority over every other method. The proposal must earn support through its outcomes.\par

% R152 source proposal; provenance and exact prior text remain in the proposal manifest.
\section{A grounded conditional learner with distinct output interfaces}
The retention question now needs a mechanism: what information would guide
the movement, and what independent observation could show that the guidance
was useful? Follow one action through the proposed measurements and learner
before adding a writer or attributing a gain to biological teaching.

Start with a movement we can measure well. For the proposed snap showcase,
the arm begins extended with the palm down. The elbow and forearm swing and
turn; the wrist and hand turn upward during the snap, then move down to rest.
The model must account for that whole movement, including the preparation,
contact and return. Finger pairing, organism-specific anatomy and timing still
need a real-reference series; a human-looking animation cannot establish that
a particular monkey can perform the same gesture.

The cup reach supplies a simpler precision comparison: relate recorded signals
to measured movement before attempting the coordinated snap. Each task still
needs its own anatomy, contact, timing and independent comparison.

That comparison connects three interacting groups: task-linked brain
observations, named-muscle observations and a musculoskeletal simulation.
A learner would relate the observations to the mechanics; independent movement
and contact measurements would test its predictions. EEG records scalp
electrical activity and surface EMG records muscle electrical activity.
A separately recorded nerve or spinal signal needs its own declared channel.
The proposed apparatus uses two arm sleeves, scalp EEG and synchronized,
calibrated left and right tripod views. Vendor, channel, bandwidth and clock
qualification remain open.

Intent to action is the target. By understanding and measuring the intermediary,
we aim to make that mapping as close to one-to-one as possible within a declared
task. Specify what counts as the intended action, the participant and session,
state and context, permitted history, tolerated error and independent reference.
The test then asks whether observing the middle reduces ambiguity and improves
prediction or qualified control. It does not assume a unique hidden intention
or claim that the best possible mapping has already been achieved.

The learner's first test would compare its predictions with recorded brain
and muscle activity and with measurements of where the hand moved and when.
Sensory feedback and physical action would add a further question: what did
each separate output interface actually deliver? A successful prediction
cannot answer that question, and EEG and EMG recorders alone cannot write
sensation. For this proposal, no paired trial, trained model, delivered sensory
output or physical actuation is established. Sources and permissions, physical
and computational behaviour, clinical interfaces, reference comparisons and
held-out tests still require independent assessment.

The retained LINK comparison belongs to a different experiment. Its 18 runs
and 720 saved epochs are not paired EEG/EMG/movement trials; both model-admission
gates and all five benefit claims failed. The prospective methods below would
have to supply the missing paired evidence and separately test recording,
prediction, simulated movement, sensory delivery and physical action. A result
for one of those jobs would not establish the others.

\subsection{Form, function and action need a measured middle}
A structural map constrains what a system could do. A NeuroConnectome maps
specified neural connections; its organism, coverage, resolution and source
would have to be declared. The signals passing through those connections
still have to be observed or inferred, with the difference made explicit.
To predict the next action, we need the processes that turn structural
constraints into an outcome. The proposed programme follows form to function
to action, while intent to action names the task-level translation we want
to improve.

DNA to RNA and RNA to protein put a process between a starting constraint and
its downstream product. Regulatory context, RNA processing and cell state
change what that process produces; protein structure alone does not determine
the behaviour of a whole cell. AlphaGenome, RNA or cell-state models and
AlphaFold illustrate these different modelling jobs. AlphaTranscriptome names
a desired predictive capability here. Connecting the predictors would still
leave system dynamics and the environment to account for, with experimental
outcomes needed to test the combined prediction.
% Source-role pointers: alphagenome2026, scgpt, alphafold3; bounded owner checks.
\par\noindent \href{https://www.nature.com/articles/s41586-025-10014-0}{AlphaGenome regulatory prediction}~\cite{alphagenomeRetained};
\href{https://www.nature.com/articles/s41592-024-02201-0}{scGPT cell-state modelling}~\cite{scgptRetained};
\href{https://www.nature.com/articles/s41586-024-07487-w}{AlphaFold 3 structure prediction}~\cite{abramson2024}.
These provide modelling examples, not evidence of an integrated cell programmer.
Our account of AlphaGenome and scGPT is limited to their reported modelling roles;
their complete methods have not been independently assessed here.

In the biological pathway, Brain2Spine and Spine2Body name proposed links to
test; their composition is the Brain2Body aim. The spinal cord integrates
descending and sensory signals and contributes local coordination, reflex
regulation and adaptation. This is a different mechanism from RNA processing,
but it serves the same explanatory purpose: account for the context-dependent
process between an initial constraint and its outcome. Supported spinal
observations would test that account. Where they are absent, the internal
states remain inferred or simulated. Physical spinal work remains read-only.
\par\noindent \href{https://pubmed.ncbi.nlm.nih.gov/38603479/}{Spinal sensorimotor adaptation}~\cite{lavaud2024}
is a constrained example, not evidence of unrestricted skill transfer or
validation of our proposed pathway. Our comparison draws on the abstract and
relevant author-manuscript passages; equivalence to the final published edition
and the complete methods and statistics remain unverified here.

Thought2Action relates a declared intention to an independently measured action.
Action2World follows the action to its consequence in the task environment;
Thought2World names the proposed composition. Body2World drew attention to the
consequence outside the body; Body2Action focuses on body mechanics and a
specified task outcome, while Brain2Motion focuses on the movement pathway.
These working labels identify related scopes for a test. They do not establish
interchangeable or working controllers.

The engineering route connects neural activity to musculoskeletal simulation
and then to a programmable model, toward NeuroTranslation. The proposed
movement-focused learner is a large movement model. System 3 thinking names
the proposed human--AI supervisory partnership around it. The supervisory
model would compare task strategies and diagnose where neural or mechanical
predictions miss independent measurements, then propose training or calibration
comparisons for the signal mapping and virtual body. The task learner would
perform the declared prediction or control; the supervisor would help choose
what to test. Neither role replaces the biological spinal intermediary or a
physical interface. Final held-out outcomes must remain outside rewards and
model selection.
\par\noindent Shaw and Nave's SSRN preprint,
\emph{Thinking---Fast, Slow, and Artificial}, and its
\href{https://executiveeducation.wharton.upenn.edu/thought-leadership/wharton-at-work/2026/05/thinking-fast-slow-and-artificially/}{institutional discussion}
provide conceptual context for this proposed human--AI role.~\cite{shawNaveSystem3Retained}
The comparison is limited to that framing; the complete preprint and institutional
article have not been reviewed here. Neither establishes a third anatomical brain
system or validates this programme's proposed learner or supervisor. Expert footage may supply
movement and tactical examples, with rights and camera uncertainty accounted
for; it supplies no measured expert EEG, EMG, muscle force or stimulation plan.

Action changes the outcome. Its sensory consequences inform perception and
the next intention. Independently measured movement would anchor the action
part of that loop; it would not by itself decode a unique sensation or
perception. Known stimuli, measured feedback and appropriate recordings or
reports would have to test those further mappings. The return path must
therefore be measured rather than drawn into existence. The full
Interoceptomics aim includes internal-body sensing and connectivity; external
touch and proprioceptive tasks contribute evidence but do not by themselves
solve that wider problem. A prediction, assisted execution and retained skill
without assistance answer different questions. Each wider claim must earn
its own independent comparison.

\subsection{Context connects. Grounding checks}
Within a reach, context has a concrete job: relate brain activity to muscle
activity in this task and at this time. Attention is one proposed way to
represent those relationships. The model compares a query with keys and uses the resulting weights to combine their values. Masks keep unavailable observations and future
samples out. The calculation weights information; it does not establish that
attention understands all context.
% Source registry IDs: context-native-01a12f62, context-attention2017-v7, context-mult2019, context-policy4
\par\noindent \href{https://arxiv.org/html/1706.03762v7}{Vaswani et al. (2017)}~\cite{vaswani2017}; \href{https://aclanthology.org/P19-1656/}{Tsai et al. (2019)}~\cite{tsai2019}. These methods motivate the computation; validation of the proposed paired system still requires the independent observations below.

A pattern can fit inside the model and still describe the wrong movement.
Grounding therefore means testing the declared prediction against qualified
observations under the same conditions. Independently measured movement would
check the outcome, while muscle activity would add an observation that
constrains the brain model. Even a matching reach cannot identify a unique
private experience. No acquired paired data or trained model is established
for this proposed loop; those comparisons still have to be made.

\subsection{Muscle activity tests the brain model}
Record brain and named-muscle activity during the same reach. In the proposed
paired trial, EEG would sample scalp electrical activity and surface EMG would
sample activity associated with the named muscles. Paired training could use
the muscle observation to constrain their relationship. Then withhold the
muscle target and movement reference from the EEG-only predictor: can the
permitted brain observations predict a specified muscle-activity observable
or movement?

Pairing and calibration remain proposed. Shared cues and movement artefacts
could give that predictor a shortcut, so the comparison needs controls for
both. The muscle observations would cover only the named muscles, and EMG
supplies neither afferent nerve activity nor direct force. Sensory claims
still need known stimuli and measured feedback. Several neural states could
remain compatible with even a well-matched pair; the pair could narrow a
task-specific explanation without revealing the complete hidden brain.
Checking the movement itself requires the next independent observation.
% Source registry IDs: context-native-01a12f62, context-policy4, overhaul-brantley2018
\par\noindent The candidate reference at \url{https://www.nature.com/articles/sdata2018133} still needs a check of its exact support and reuse rights.

\subsection{Measured motion tests the virtual body}
Measure the same hand and arm from the aligned left and right tripod views.
The proposed motion-capture relay-composite would align the two point maps
in a common clock and coordinate frame, then join them into one composite
while retaining each point's origin. Before that joined map can check a
prediction, its calibration, alignment and capture precision would have to
be independently verified.

The movement map would constrain the external movement model: surface motion
and pose in the task's physical frame. To explain that motion mechanically,
the internal musculoskeletal model would describe specified bones, joints
and soft tissues. Organism, developmental stage and task would determine the
appropriate anatomical sources. X-ray projections could inform bone and
joint relationships; appropriate MRI references could inform specified soft
tissues. Registration, rights and source-specific limits would need independent
checks. No scan set is selected for this proposed comparison. Anatomy still
leaves activation, forces and neural function to test.

With a qualified composite, the musculoskeletal-simulation physics-engine's
joints, muscles, contacts and physical limits would face a measured comparison.
Fit the model on training data, then compare its predicted reach with
held-out measured movement. A force claim would require an additional force
measurement. The body's form limits possible motion; measured function would
test the mechanics and show where they need refinement.
% Source registry IDs: context-native-01a12f62, context-component-names-01a12f81, context-policy4, context-opensim2018, context-myosuite-docs, context-precision-01a12f77
\par\noindent \href{https://journals.plos.org/ploscompbiol/article?id=10.1371/journal.pcbi.1006223}{Seth et al. (2018)}~\cite{opensim2018}; \href{https://myosuite.readthedocs.io/en/latest/}{MyoSuite documentation}~\cite{myosuiteDocsUndated}. The simulation tools supply methods for that comparison; the proposed capture and body models still need its independent validation.

A lifelike hand makes that comparison visible. It cannot establish biological
fidelity. Capture alignment, precision, mechanics and biological fidelity have
not been demonstrated by this proposal, and captured geometry cannot establish
hidden neural state, internal forces or unique subjective experience. The
measured reach would check a movement prediction within those limits.

\subsection{Two pairs. One conditional learner}
The two comparisons would meet in the proposed large movement model: EEG
observations paired with muscle activity, and measured movement compared
with predicted mechanics. Large movement model names the proposed family;
deep neural network describes a possible implementation. Large refers to
the intended scope, rather than a demonstrated parameter count. Encoders
would receive only the observations permitted for the named test. Two
branches would summarize the relationships and meet in a shared task
representation; separate output parts would predict each direction's targets.
% Source registry IDs: context-native-01a12f62, context-policy4, context-attention2017-v7, context-mult2019
\par\noindent \href{https://arxiv.org/html/1706.03762v7}{Vaswani et al. (2017)}~\cite{vaswani2017}; \href{https://aclanthology.org/P19-1656/}{Tsai et al. (2019)}~\cite{tsai2019}. These computational methods offer precedents for a design, leaving the proposed paired learner and its performance to be tested.

Signal features belong to branches, and branches belong to the learner.
That hierarchy describes the computation, not two anatomical thinking systems.
A common clock, task and physical frame would connect the branches; independent
measurements would check their errors. Form constrains the prediction, and
checked function would guide refinement. Several hidden states could still
fit the same reach, so neither a shared representation nor conditional
predictions would establish an exact inverse nervous system. The architecture
has not been selected, implemented, trained or independently accepted.
To choose it, we would have to compare it with simpler baselines on
appropriately reserved observations.

\subsection{Learn the movement, then test what transfers}

A reach can provide a training target even when no EEG accompanies it. Earlier movement supplies the input; a later observation supplies the target. A separate proposed comparison would ask whether learning that relationship in one person helps forecast another person's grasp-and-lift movement. The WAY-EEG-GAL descriptor provides a relevant continuous-recording candidate, alongside its distinct windowed trials and derived lift summaries.\citep{luciw2014way} This movement-only question needs no AllLifts Run/trial join or invented cross-person neural pairing. The original condition-specific five-weight study remains a separate design.

The mapping has to explain what two bodies share. Homologous structures share evolutionary anatomical origin; mechanically analogous structures can serve comparable roles without that origin. A human digit aperture and a robot gripper opening could support a task-space comparison after separate mechanical qualification; this human-only proposal would not test that robotic extension. Joint degrees of freedom, axes and signs, proportions, muscle--tendon roles, loads, inertia, contact and time scales remain body-specific. Similar geometry cannot establish equivalent dynamics or neural coding. Retargeted motion remains derived; tracking, estimation, prediction and simulation keep their different evidential roles.

For the proposed two-person comparison, the representation contains index-, thumb- and object-sensor positions relative to the wrist, digit aperture and causal wrist velocity: thirteen features, with thirteen availability masks. Each participant needs qualified position units, a proper wrist-frame rotation, orientation convention and sensor-attachment calibration. A positive median wrist-to-index-sensor distance, expressed in qualified metres, supplies the characteristic length from assigned training support only; it is not a bone measurement. Multiplying normalized velocity by the same fixed 0.1-second reference interval makes all thirteen features dimensionless. Physical velocity errors invert that scaling with length divided by 0.1 seconds. Missing features receive no target loss; an unavailable observation is not a measured zero. Full finger joints, individual muscle forces and subject-specific inertia are not supplied by these sensor features.

The continuous target would be the same observed feature set 100 milliseconds beyond the latest actually available filtered input time. The proposal retains fifty frames at 100 Hz and a causal 51-tap filter before decimating qualified 500-Hz kinematics by five. Its 50-millisecond linear-phase delay must remain in the physical timing account; shifting information backward would change the prediction claim. Coefficients, passband, timestamp convention, gaps and boundary handling still require qualification. Only named kinematic columns may enter: the documented kinematics matrix also contains force-plate signals.

Identity and independence answer different questions. Trustworthy versioned participant-release attribution, an immutable selected archive entry and a qualified capture can identify the participant record and input object without requiring one particular internal ID-field spelling. They cannot establish biological identity or independent original acquisitions. Shared samples, trims, derivatives and overlapping recordings must form indivisible acquisition groups before filtering, windows, normalization or role assignment. Different names, checksums or series numbers do not prove separation. The complete public non-ST candidate roster and these relations still need establishing; all 294 rows of the prior AllLifts table remain excluded from training, evaluation and assigning trials or acquisition groups, and ST, competition storage and the reserved LINK recordings remain unavailable.

Qualified P1 groups would supply source pretraining; P2 would supply target training and development. With $G\geq2$ independently established P2 acquisition groups, the fixed identifier-hash ranking assigns the first $\lfloor3G/5\rfloor$ to training and the remainder to development. Canonical original-parent identifiers, their serialization and the roster must be fixed before ranking or examining outcomes. This is a nominal 60/40 split, not chronological ordering or a promise of condition balance. If unresolved ancestry merges P2 into one group, the comparison stops; splitting its windows cannot manufacture an independent development group. No group or segment has been assigned by the proposed identity method.

Three matched arms would test the transfer. B0 pretrains on correctly mapped P2 training movement. B1 and B2 use identical P1 windows and order: B1 sets the index-role input to the thumb sensor, the thumb-role input to the object sensor and the object-role input to the index sensor, while B2 retains the correct correspondence. Both share the correct source calibration and scale before that corruption, keep the wrist anchor, and use the same correct target adapter and ordered target batches. Parent-balanced sampling must record unique windows and repetitions. The false mapping can be relearned or can alter feature availability; matched update counts do not match source diversity. These are limits on attributing any difference to anatomical correspondence, not reasons to replace the control after seeing its result.

The proposed learner is one GRU layer of hidden width 32, with 26 inputs and thirteen outputs. Each arm and fixed seed---17, 31 and 73---receives 200 pretraining and 100 target updates, batch size 16, identical per-seed initialization and the fixed last checkpoint. Adam uses learning rate 0.001, betas 0.9 and 0.999, epsilon $10^{-8}$ and zero weight decay, without clipping, scheduling or early stopping. Training minimizes mean squared residuals over observed target features; a batch with no observed targets is rejected. The resulting 2,700 updates are planned, not performed. No external checkpoint is selected; numerical updates, filtering, precision and actual workload limits remain to be checked independently.

For each development parent and seed, the proposed score takes the square root of summed squared errors divided by the count of observed target features, then averages parent scores without weighting by parent duration. A zero denominator rejects the comparison. Report feature counts, missingness, physical-unit errors, uncertainty and persistence at the same information endpoint, including every worse parent or feature. B2 would have to lower normalized error against both B0 and B1 in every fixed seed, with no applicable original criterion failing, to justify another investigation. This remains a two-person development comparison, not independent final confirmation, population transfer or biological teaching.

Broader motion collections could supply additional movement examples only after their exact assets, rights, task coverage and original-recording overlap are qualified. In \href{https://www.addbiomechanics.org/download_data.html}{AddBiomechanics}, measured markers and ground reactions must remain distinct from fitted kinematics and estimated inverse-dynamics torques. \href{https://amass.is.tue.mpg.de/license.html}{AMASS} requires checking the selected asset's research-use, non-commercial and redistribution restrictions. Copies across collections stay with their original acquisition. LaBraM's WAY source-family exposure is flagged, while exact checkpoint/participant/trial exposure remains unknown; reserving a recording later cannot erase earlier pretraining exposure. No such checkpoint is selected for this comparison.

Before adding X-ray/CT or MRI, inspect the anatomical sources already used by the selected musculoskeletal model. A scan is useful when it constrains a stated gap or supplies an independent check beyond that model. Its person or species, scale, registration, rights and uncertainty still matter. Another person's scan remains an external anatomical reference; a generic model does not implicitly measure this participant's geometry, activation or forces. No scan set is selected here.

Actual source eligibility, acquisition groups, units, frames, clock and calibration remain unresolved for this continuous comparison, so no new fit or successful transfer is reported. A useful movement prior could later be tested with actual paired neural/body observations and the permitted input masks. Its transferred targets would still not become simultaneous measurements. The next result must identify where the mechanical forecast improves, where it worsens and which measurement or assumption accounts for the difference; neural equivalence, causal physiology, sensory delivery and retained skill would require further experiments.

\subsection{Pretrain. Pair. Verify. Refine}
Those comparisons start with the observations we may use. Data rights,
calibration and held-out subjects, sessions and tasks determine which can
enter the study; resolve them before processing data or choosing a model.
Open-source code does not grant blanket permission to use datasets,
checkpoints or clinical interfaces. Qualified data could first train
representations for each recording type. LaBraM is one EEG encoder candidate,
not a complete brain-body writer. Synchronized paired observations could
then train the declared mapping from permitted inputs.

Paired prediction would give later refinement a baseline to beat. Before
optional reinforcement learning, post-training or task-specific fine-tuning,
the physical predictions and reward need checking against independent
outcomes. Offline reinforcement learning needs qualified coverage of actions,
transitions and rewards: labels alone cannot tell us what an untried action
would do. Backpropagation can update differentiable neural components, but
it cannot automatically pass through every physics engine; universal
differentiability is not established. Simpler baselines and frozen final
tests remain through every phase.
% Source registry IDs: context-native-01a12f62, context-policy4, context-labram2405, context-cql2020-abstract, context-opensim2018, context-myosuite-docs
\par\noindent \href{https://arxiv.org/html/2405.18765v1}{Jiang, Zhao and Lu (2024)}~\cite{labram2024}; \href{https://arxiv.org/abs/2006.04779v3}{Kumar et al. (2020)}~\cite{cql2020}; \href{https://journals.plos.org/ploscompbiol/article?id=10.1371/journal.pcbi.1006223}{Seth et al. (2018)}~\cite{opensim2018}; \href{https://myosuite.readthedocs.io/en/latest/}{MyoSuite documentation}~\cite{myosuiteDocsUndated}. These sources offer candidate methods for the sequence; validation of the proposed paired system still requires the experiment below.

No data acquisition, training, fine-tuning, RL rollout or accepted checkpoint
is established for this proposed loop. The next experiment must test what
the prediction and reward represent against independent outcomes. A reward
that the computer can calculate supplies no biological proof on its own.

\subsection{Respond quickly. Improve carefully}
\label{b2-frozen-response}
The proposed response path would use a frozen accepted version to estimate
the current state, predict declared targets and prepare encoded outputs for
a separately qualified interface. Permitted causal time windows would limit
it to the allowed past and present observations. Estimating the state of this
moment would not change learned weights. A separate adaptation path would
inspect errors, calibrate or recalibrate, and propose a candidate update;
independent evaluation, a restorable checkpoint, restoration proof and
applicable authority would precede promotion. Examining errors would not
silently change the running model, and no unchecked live weight promotion
would occur.

The fast-and-slow comparison is a Kahneman-like analogy for this proposed
division of work. It identifies neither two biological or anatomical systems
nor inherent thinking speeds for EEG, spatial data or other sensor modalities.
Real-time performance would have to be measured across acquisition, filtering,
alignment, inference, output preparation, transduction or actuation, feedback
and rendering. Report delay tails, jitter, missed deadlines, uncertainty and
coexistence with healthy workloads; component inference time alone cannot
establish end-to-end performance. No accepted model, running system that adapts
under these rules or real-time demonstration is established here. The loop has
no measured end-to-end latency or accuracy guarantee.
% Source registry IDs: context-native-01a12f62, context-policy4

\subsection{Write sensation. Enact action}
\label{b2-output-interfaces}
The prediction ends at the computer. A separate interface would have to turn
it into sensation or action. Input and output describe direction relative to
that computer; afferent and efferent describe physiological direction relative
to the central nervous system. A computer output could drive a sensory
interface supplying input toward the nervous system, while a computer input
could record activity related to the body's motor output. Both boundaries
must remain explicit.

The sensory route would use encoders and a separately qualified visual, haptic
or other sensory transducer to deliver specified feedback. Direct neural
writing remains a separate later research interface, with its own device,
measurement, safety and clinical evidence. Direct cortical writing could have
a sensory role without being a peripheral afferent signal. Predicting an
EEG-like pattern would not deliver sensation: the sensory content, interface,
authority and actual measured response would need their own tests, including
conditions, error and delay.

The action route would send separately qualified encoded controls through a
controller to a virtual body, a robotic actuator or a later qualified biological
interface. Action encoders, controllers and actuators have distinct jobs in
that route. A robot command is an action-role output, not an anatomical efferent
nerve. Simulated action, physical actuation, retained learning and biological
causal claims each require their corresponding evidence. EEG and EMG recorders
neither supply sensation nor actuate a body. These sensory and action
subinterfaces belong within the unchanged five programme pillars.

Failed access to a clinical source supplies no methods evidence. Implanted
precedents cannot establish that EEG or EMG recorders are equivalent to devices
that deliver sensation or produce an action. The acquisition and virtual-modelling
plan excludes physical sensory writers, transducers and body actuators; these
additions remain unqualified and unpriced. Acquisition and virtual-modelling
costs do not establish what delivering sensation or action would cost.

This recording-and-virtual-modelling proposal cannot test physical sensory
delivery or biological actuation. No such physical interface, stimulation,
acquisition or actuation is established here, and no loop, training or
physical-writing result is claimed.

A physical study would need its own qualified interface and procedure. This
proposal supplies no stimulation dose, implantation, electrode-placement,
acquisition or actuation procedure. It authorizes no stimulation, participant
acquisition, biological actuation, new equipment purchase or training, and
grants no participant, clinical, device, purchase or deployment permission.
% Source registry IDs: context-write-native-01a12f66, context-native-01a12f62, context-policy4

\subsection{Close the loop with evidence}
Return to the cup and test each prediction against the observation it was
asked to explain. The brain-to-body prediction would face withheld observations;
the body-to-recorded-neural estimate would face its own recorded target.
A sensory or action interface would need a further test of what it delivered,
under which conditions, with what error and delay. Two models can agree and
still be wrong; a high reward can have the same problem. Neither agreement nor
a successful prediction would, by itself, identify a unique inverse brain or
establish someone's private experience.

The programme's distinction remains that approximation gets us towards
plausibility and virtualization gets us to precision. Independent measurements,
repeatability, causal checks and stated error bounds must earn that precision.
The original held-out and real-reference controls would test the prediction;
retained-learning controls would test whether a gain remains after useful
guidance ends. Precision is not an automatic result, and the proposal establishes
no perfect biology.

The next experiment must report where each declared prediction works and where
it fails against its independent measured target. Its answer would determine
which part of the proposed loop needs refinement.
% Source registry IDs: context-native-01a12f62, context-write-native-01a12f66, context-precision-01a12f77, context-policy4, context-spec15, context-glossary13, b2b-circular-checklist-original

\subsection{Technical representation and mask}
The information boundary must survive the calculation. A model given tomorrow's
sample or today's withheld target could make a persuasive prediction of an
answer it was never meant to receive. The representation and mask specify
which comparison the proposed attention computation would actually perform.

Learned query, key and value representations give the proposed attention
calculation a precise form. An admitted mask would restrict which observations
may contribute; biological grounding still requires the independent comparison.
\[\operatorname{Attention}(Q,K,V)=\operatorname{softmax}\!\left(QK^{\mathsf{T}}/\sqrt{d_k}+M\right)V.\]
Here Q, K and V are queries, keys and values; M expresses the admitted mask. The dimension $d_k$ is the common feature width of each query and key vector. The equation is methodological context from the supplied technical proposal. Modality, channel, task and time are part of the observation contract. No weights, EEG encoder or gradient route are selected or executed.

\subsection{Permitted observations and causal modality routes}
Declare B=EEG, E=EMG, Y=independently measured movement, C=task/stimulus/context, and P=the source-bound physics model. Training may relate qualified paired observations. Inference uses only its explicitly admitted subset:

\begin{itemize}
\item EEG-only: permitted past/present B and declared context predict a specified   EMG/movement target. Withheld E, Y and future samples cannot enter the predictor.
\item Body-only: permitted E and/or Y predict specified recorded neural features as   a conditional distribution with uncertainty. This is not a unique inverse brain.
\item Hybrid: additional modalities require a separately named route and evaluation;   a hybrid result cannot be reported as EEG-only.
\item Simulation-derived inputs stay simulated. They do not replace independent   observations or inherit the biological validity of a rendered body.
\end{itemize}

Separate training signals from evaluation references. Cycle consistency can make two models agree while both are wrong. Shared cues, movement artefacts, subject identity, session overlap and preprocessing leakage need the original controls. Do not expose final held-out targets to model, physics or reward selection.

\subsection{Representation, supervised learning and policy refinement}
\begin{enumerate}
\setcounter{enumi}{-1}
\item Establish source rights, subjects/tasks, measurement units, calibration,    uncertainty, clock alignment and frozen subject/session/task splits. Audit    pretraining overlap before preprocessing and model selection.
\item Pretrain modality representations on the qualified inventory. LaBraM is one    candidate EEG encoder: neural-tokenizer training, masked neural-code prediction    and downstream fine-tuning. It is not a complete brain-body writer. Code,    checkpoint, dataset and individual-use rights remain separate qualifications.
\item Learn paired supervised mappings with named targets and explicit modality    masks. Compare appropriate simpler, cue-only and artefact baselines. Keep raw,    processed and model-derived values distinct.
\item Qualify physical predictions and the reward verifier against independent    measurements. A successful task, smaller calibrated movement error and valid    physical states are prospective reward ingredients, not completed results.    Declare loss/reward weights on training/development evidence. A computable    objective or high reward does not confer biological ground truth.
\item Where admitted action/transition/reward coverage supports it, evaluate optional    RL and post-training. Offline trajectory labels alone do not establish the    consequences of actions never supported by the data. Conservative Q-Learning    is a methodological candidate, not a safety or performance guarantee here.
\item Fine-tune for the declared task and qualify both directions on frozen final    tests, uncertainty, baselines and every original threshold. Keep failures.
\item Serve a frozen accepted version. Propose later updates separately and promote    only after independent acceptance, restoration proof and applicable authority.
\end{enumerate}

An illustrative supervised objective would combine declared errors for
specified EMG, movement and recorded-neural targets with applicable physical
constraints and optional consistency terms. Each term needs its units,
measurement uncertainty and target definition before a weight is chosen.
Consistency can check whether the outputs fit together; an independent
observation checks whether they fit the measured outcome.

Backpropagation changes differentiable components of the computer model;
biological feedback is a different process. A physics engine in the path does
not automatically give the learner a gradient. Qualify the actual
OpenSim/MyoSuite rollout and optimization route, or a separately validated
surrogate, rather than assuming gradients through every route. No training,
RL rollout or task checkpoint is created or certified here.


% Proposed conceptual-article insertion only; existing manuscript untouched.
% Inherit Noto Sans/Fira Math from the selected master.
\section{A larger model still owes us a learning experiment}

A model could predict a hand's path more accurately while teaching its owner nothing. That distinction matters as the programme expands from recording-specific mappings toward a proposed Large Movement Model. The new hypothesis is that greater amounts and greater diversity of qualified neural, muscle and movement data could improve generalization. Brain2Body asks what any resulting prediction would let a trainee do, what information or assistance supplied the gain, and what remains when teaching ends. A scaling result would answer the first modelling question only within its tested domain.

\subsection{Four training sources, five component roles}

The expanded training proposal has at least four source groups: neural observations; muscle and other myographic observations; measured motion and displacement; and data generated by physics-based motion or musculoskeletal simulation. Suitable neural sources may span brain, spinal, peripheral and relevant visceral recordings. Their actual modalities, sites and task roles remain distinct. The present emphasis is motor--sensory. Myoception is the programme's term for the sensory side of musculoskeletal contraction and movement, including proprioception, tactile feedback and other task-relevant sensations linked to action. It marks proposed scope, without establishing a new external scientific standard. That scope supplies neither acquired coverage nor body-wide simultaneous observations, and the four groups are not an exhaustive taxonomy.

Known tactile stimuli, proprioceptive feedback, contact, force and task context contribute different evidence. Stimulus delivery, reported sensation, recorded sensory activity, calibrated force and tracked configuration cannot be substituted without a tested observation relationship. EMG records muscle electrical activity, not direct peripheral nerve or sensory activity; kinematics alone does not measure proprioceptive experience. Each measured, reported, estimated or simulated quantity keeps its provenance.

The simulation group supplies potential training examples as well as constraints: trajectories, state transitions, muscle--tendon or contact-force outputs and declared movement priors. Record the simulator, body-model version, parameters, solver, control representation and underlying sources. Joint-target control must not be called muscle-driven control. Retargeted recordings retain original lineage and transformations; generated states or forces remain simulated. The generator requires calibration against appropriate measured support with uncertainty and training-only role boundaries.

Test its contribution by comparing a fixed-real-data reference with a simulation-augmented arm under matched permitted information, compute and original criteria. Record real-example exposure, simulated exposure and updates in each arm; a fixed real corpus alone does not match training exposure. Count simulated amount and diversity separately from real acquisitions. Held-out real observations must check the result without entering generator calibration, tuning, reward choice or model selection. Wrong dynamics, contacts or source correspondence can make the synthetic contribution fail. Generated data is never independent biological ground truth for the mechanism that produced it. Generator methods, implementation and use conditions remain to be checked; no frontier implementation or dataset is selected here.

Five component roles describe another axis: brain signals, muscle signals, measured movement, a virtual body and a learning model. A virtual body can supply training data and physical constraints, while a learning model can draw from several source groups. Neither count excludes another useful, properly characterized observation. These resources give the later teaching study more hypotheses to test, rather than supplying the missing retained-learning result.

\subsection{Breadth needs a controlled comparison}

Amount and diversity must be varied separately. Repeated windows from one acquisition increase presentations without adding independent subjects, days or tasks. A proposed amount comparison would increase support from a stated source distribution. A diversity comparison would retain a declared amount while changing named aspects of that distribution, with a crossed comparison testing their interaction. Organism, subject, session, movement, body region, modality, apparatus, environment and measurement quality require separate descriptions. Dataset count alone cannot measure them.

The architecture, permitted input information, targets, preprocessing and selection rules would remain fixed. Source rosters, repetitions, updates and measured computing work must accompany the results; equal update counts alone do not establish matched compute. If diversity also changes recording quality or task difficulty, the comparison cannot attribute its outcome to diversity alone. All exact budgets, controls, seeds, uncertainty methods and decisions remain to be registered before a new experiment. This proposal changes neither the adopted continuous WAY-EEG-GAL development design nor any original protected partition.

A candidate amount--error curve, for a fixed diversity condition and evaluation domain, is
\[
E(N)=E_{\infty}+A(N/N_0)^{-\alpha}.
\]
$N>0$ and its fixed reference $N_0>0$ use the same declared training-amount units. $E$ is a chosen nonnegative error; $E_{\infty}$ and $A$ use its units, while $\alpha$ is dimensionless. This decreasing-error candidate requires $E_{\infty}\geq0$, $A\geq0$ and $\alpha>0$: $A>0$ gives strict decrease, while $A=0$ gives constant error. Positive $\alpha$ alone is insufficient. Predictions must also respect any upper bound of the metric throughout the declared fitting and prediction range. The floor is a prospective curve parameter, not an established irreducible biological limit. No exponent is estimated here and no universal law follows. Constant, saturating, increasing, non-monotonic and domain-specific alternatives need prospective comparison and predictive checks at withheld scaling levels. Observed negative transfer remains a result even when this constrained candidate cannot fit it. Related windows cannot supply independent uncertainty replicates. Neural-target, EMG-target, pixel and physical-motion errors retain their own units and denominators.

Generalization would be evaluated separately across excluded subjects, sessions, tasks, datasets and, where mechanically qualified, species. The target-only/no-transfer reference, applicable incorrect correspondence, simpler predictors and original null, shuffle, cue/confound and causal-time controls remain. Negative transfer, worse domains and saturation stay in the record. A better pooled score cannot discharge an original per-output no-harm criterion or erase known pretraining exposure.

\subsection{Build a composite without inventing a shared observation}

FALCON M1-A and AJILE12 illustrate complementary source candidates. The fixed FALCON M1-A release describes simultaneous M1 activity and processed intramuscular EMG from one rhesus macaque; released synchronized motion capture is not established. Its acausally processed EMG targets limit causal real-time interpretation. The small minival subset is drawn from calibration data for submission-format checks and is not independent validation. AJILE12 describes intracranial recordings and two-dimensional, video-estimated upper-body pose from twelve clinical participants. Paired EMG and instrumented three-dimensional mocap are not established there. Its pose-label construction is participant-adapted; particular held-out-day estimator exposure and numerical alignment uncertainty remain unqualified. \href{https://doi.org/10.48324/dandi.000941/0.241029.1405}{FALCON M1-A release}; \href{https://papers.nips.cc/paper_files/paper/2024/file/8c2e6bb15be1894b8fb4e0f9bcad1739-Paper-Datasets_and_Benchmarks_Track.pdf}{FALCON methods}; \href{https://doi.org/10.1038/s41597-022-01280-y}{Peterson et al., 2022}; \href{https://doi.org/10.48324/dandi.000055/0.220127.0436}{AJILE12 release}.

Qualified source-specific pretraining could exploit each relationship. A transfer model would then have to demonstrate a correspondence on actual target measurements. The separate recordings never become one simultaneous neural--muscle--movement event. The release descriptions do not establish downloaded recordings or a trained model. Their exact use conditions, versioned files, integrity, original acquisition relationships, timing, units, earlier exposure and independent partitions still need checking before training. Retargeting between homologous or mechanically analogous structures also requires declared geometry, frames, scale, load, contact, timing and uncertainty. Similar movement cannot establish identical neural coding.

\subsection{Model the inside and the outside of one action}

When a hand closes on an object, external displacement and contact put demands on internal mechanics. The proposed world model would represent specified musculoskeletal states alongside the supported body, objects and environment. NeuroTranslate would connect named observations and targets; MoCap ReComp would provide its separately qualified source-preserving movement composite; MyoSim Replay would simulate declared mechanics. Their proposed combination is the Large Movement Model. Neither that name nor a neural composite implies a measured whole nervous system, equivalent cells across organisms, complete internal organs or unique access to private experience.

The model must distinguish measured neural activity, muscle electrical activity and tracked motion from estimated activation, internal forces, latent neural state and simulated motion. EMG is not direct peripheral nerve recording or automatically calibrated force. A proposed transition model could forecast a distribution over joint configuration, generalized velocity, muscle-model activation and environmental state from permitted past observations, context and declared controls. Separate observation models would connect each sensor to the state it can constrain. An unavailable observation stays unavailable; a conditional reverse prediction cannot identify a unique hidden brain.

Physical tests would compare aligned predicted movement with independent measured motion, and any internal-force claim with appropriate independent force evidence. Joint limits, contacts, loads, energy accounting, friction, gravity and numerical solver behaviour need their own checks. OpenSim supplies established methods for declared musculoskeletal models; it cannot confer biological validity on this proposed assembly. \href{https://journals.plos.org/ploscompbiol/article?id=10.1371/journal.pcbi.1006223}{Seth et al., 2018}. A rendering can show where a model says a muscle pulls. Only the applicable observations can test whether that pull explains the claimed result.

% R4 conceptual-article insertion candidate; no manuscript/source promotion.
% Inherit the selected master's Noto Sans/Noto Sans Mono/Fira Math setup.
\subsection{One muscle pulse, two simulated trajectories}

A moving virtual hand has answered a mechanical question before it has answered a teaching question. In one numerical diagnostic, we compared two trajectories of the installed \texttt{myohand\_r} human right-hand mechanism on a fixed passive torso. A muscle excitation was added to one branch. The other supplied its matched control. This asks whether the specified software input changes the computed state. It provides no neural decoding, measured biological response or retained-learning outcome. The mechanism used \texttt{myo\_sim} 0.2.3, commit \texttt{93b0ca8f4ec90c9899ee7f05fee561e9911da91b}, MuJoCo 3.14.0 and NumPy 2.5.3.\cite{myohandMyoSimPinned2026,myohandMuJoCo3140,mujoco2012,myohandNumericalRecord20261105}

\subsubsection{Mechanism and matched comparison}

The compiled model has 23 scalar hinge degrees of freedom, 39 muscle--tendon actuators, 39 activation states and zero sensors. Both branches began with matching full recorded initial states, including positions, velocities and activation states. The control branch received no additional excitation pulse; the pulse branch received dimensionless excitation 0.2 at \texttt{FDS2\_r}, with all other controls zero. Excitation is a muscle-model input, distinct from a joint-position request, measured EMG or neural activity. Gravity and passive dynamics remain present, so the unpulsed trajectory need not remain stationary.

The native Euler integrator advanced each branch through 100 explicit steps of $0.002\,\mathrm{s}$, giving nominal $T=0.200\,\mathrm{s}$. Pulse steps 10--59 inclusive correspond to the pre-step interval $0.020\leq t<0.120\,\mathrm{s}$. Each branch received explicit initial and final forward evaluations to synchronize state-dependent quantities. Thus the record contains 200 explicit integration steps and four explicit forward evaluations. The executed source path also implies two compilation calls and one construction-time internal forward evaluation. Those counts do not enumerate every internal native solver operation.

The construction path changed physical parameters even with optional \texttt{inertia\_floor=None}: each diagonal inertia component of four fixed scaffold bodies changed from $1$ to $10^{-4}\,\mathrm{kg\,m^2}$. This native transformation is part of the tested mechanism, rather than evidence that all parameters remained untouched. No additional physical retuning or \texttt{implicitfast} fallback was used. Position, orientation, scale, contact behaviour and biological fidelity have not been independently calibrated by this diagnostic.

\subsubsection{Recorded response}

For pulse branch $P$, control $C$ and the synchronized final state, the reported quantities are
\[
\Delta q_{\infty}=\max_{1\leq j\leq23}|q_{P,j}(T)-q_{C,j}(T)|,
\qquad
\Delta a=a_{P,\mathrm{FDS2}}(T)-a_{C,\mathrm{FDS2}}(T),
\]
\[
\Delta F=F_{P,\mathrm{FDS2}}(T)-F_{C,\mathrm{FDS2}}(T).
\]
The first quantity is the maximum absolute separation across the 23 hinge coordinates at the final time, in radians. The other quantities are signed pulse-minus-control contrasts in the selected dimensionless activation and tendon-actuator force in newtons. Force follows the model's convention and uses synchronized endpoint forward evaluations, rather than derived values left over from the last integration step.

% Shared by all four proposed R3 inserts; no added packages or font defaults.
\begin{table}[htbp]
\centering
\caption{Final contrasts for one simulated right-hand mechanism. All values compare the pulse branch with its matched control at nominal $T=0.200\,\mathrm{s}$. The angle value is the largest absolute separation across 23 hinges at that endpoint; activation and force retain the pulse-minus-control sign. Activation and force concern the selected \texttt{FDS2\_r} actuator. Force uses the synchronized final forward evaluations and the model's sign convention. These are simulated contrasts, not errors against biological measurements. The displayed digits preserve the recorded numerical values, without establishing biological accuracy.\cite{myohandNumericalRecord20261105}}
\label{tab:myohand-numerical-response}
\begin{tabular}{p{0.53\linewidth}p{0.26\linewidth}p{0.10\linewidth}}
\hline
Quantity & Contrast & Unit \\
\hline
Largest absolute final joint-coordinate separation, $\Delta q_{\infty}$ & $0.5027347840656727$ & rad \\
Selected final activation difference, $\Delta a$ & $0.05205729904534248$ & none \\
Selected final tendon-actuator force difference, $\Delta F$ & $-6.407722136774943$ & N \\
\hline
\end{tabular}
\end{table}


Independent auditing of the saved evidence completed 132 standard-library checks without a failed check or model rerun. All required recorded numerical channels were finite and recorded native warning counts were zero. Four stored endpoint mass-matrix spectra were positive. Raw mass matrices and residual operands were not independently exported; their stored diagnostics were checked without independently resolving those matrices. No contact was observed. At the final control endpoint, the recorded soft-limit excess was $0.0004121829780273552\,\mathrm{rad}$, so exact hard-bound enforcement is not established. Neither conservation, numerical convergence nor a physiological response threshold follows.

These contrasts establish a response to the added excitation within this simulated configuration. They are not errors against measured human movement, EMG or force. Two dependent trajectories with 200 time steps supply neither 200 independent experiments nor biological replication. No statistical $p$-value is inferred, and numerical precision in the record does not establish biological accuracy or repeatability across cases.

\subsubsection{From a response to a teaching experiment}

The run used CPU dynamics. Readback showed a one-thread BLAS pool and dynamics pool size zero; transient compiler parallelism was unmeasured. Its 263 observed loaded-file identities provide bounded identity evidence, not complete system closure. Resource and performance qualification remain unresolved. The reported process high-water RSS of 5,187,464 KiB can retain pre-exec history and cannot be assigned to this model or establish a model-only 3 GiB bound. Tool/native timing near 448 ms and in-probe timing near 491.927 ms have unreconciled boundaries. There is no accepted end-to-end speedup, useful complete latency or full CPU/GPU partition.

The hand's zero sensors leave Myoception---the programme's proposed sensory side of muscle-driven action---without a measurement in this test. Simulated activation, force and motion do not supply afferent recordings, biological feedback, a learned controller, a trained Large Movement Model or generalization. The four proposed source groups remain neural observations, myography, measured motion and physics-generated examples. Mechanics, controls, calibration, uncertainty, source rights and data lineage require qualification before this generator could enter training, and its contribution would need independent held-out real tests. The amount/diversity hypothesis remains untested.

For Brain2Body, even a later qualified controller would leave the teaching question: what does a trainee retain when trainer commands end? That comparison still needs the original matched information, practice and assistance controls, sham/delayed/unrelated guidance, independent outcomes and a constant compensatory movement bridge where required. This diagnostic contains no trainee, teaching, withdrawal or physical intervention. It selects no training generator and supplies no biological ground truth. LINK's eighteen fits and 720 saved epochs retain both failed admission gates and all five failed benefit claims; protected tests and quarantined sources remain closed.

\subsection{Carry the test through teaching and withdrawal}

A better movement prior might improve assistance, but the teaching hypothesis requires another comparison. A prospective trainer--trainee study must identify the model's inputs, expert information, feedback, commands, practice and assistance burden. Match those resources against the appropriate simpler controller and original sham, delayed and unrelated-guidance controls. Otherwise an apparent neural benefit could come from supplying the target action, more practice or a less demanding task.

Retention is assessed after trainer commands stop, with any necessary compensatory movement bridge held constant. Immediate performance, delayed retention and unfamiliar task variations remain separate outcomes. A gain that depends on continued guidance may support assistance. A controlled gain that persists would need its own domain and uncertainty; it would still establish neither tissue repair nor arbitrary skill transfer. Simulation, offline replay, measured biological feedback and physical intervention retain distinct evidence and authority.

No scaling result, trained body-wide world model or biological teaching result is reported here. LINK retains eighteen fits and 720 saved epochs, with both model-admission gates and all five benefit claims failed. A broader programme gives those failures somewhere useful to lead: a new, source-qualified comparison that can show exactly which prediction improves, what supplies the improvement and whether a trainee retains anything from it.

% Prospective assistance/retention protocol; no biological experiment reported.
\section{From a guided action to a retained skill}

Catch a ball with outside guidance, then try again after the guidance ends. The
second catch asks what changed in the trainee. The comparison is useful only if
the task and access to movement remain comparable; removing a necessary movement
bridge along with the trainer would confound the test. This distinction carries
the proposed sequence from touch and object handling to NeuroBoxer and NeuroArcher.

The near-term apparatus would record scalp EEG and surface EMG from sleeves on
both arms, alongside synchronized, calibrated left and right tripod views.
OpenBCI or an equivalent configurable system remains a candidate requiring
channel, site, reference, bandwidth, unit, delay and clock qualification. Each
event needs its independently observed movement and, where claimed, an
appropriate contact or force reference. Surface EMG supplies neither direct
force nor afferent nerve activity. A virtual body can predict a response, but
its predicted contact or sensory state cannot stand in for those observations.

This left/right hand-task apparatus and the earlier opposing top/bottom dot-map
apparatus are different proposed configurations of the motion-capture
relay-composite. Their calibration, frames, visibility, occlusion and timing
must be qualified separately. Preserve the source of each measured point and
the uncertainty of the joined map. The musculoskeletal-simulation physics-engine
also needs task- and organism-specific anatomy and mechanics; lifelike rendering
cannot establish its function.

\subsection{Build coordination through four declared tasks}

\paragraph{Touch and movement.}
Rapid hand rubbing brings contact and movement together. Record the requested
knuckle movement and interlaced-hand stretch, with palms outward and arms
extended, as separate conditions. Matched observations must distinguish
contact-related effects from movement-related effects; the cold-hands metaphor
does not establish a cold-exposure result or identify a unique sensation.

\paragraph{Ball transfer and juggling.}
The proposed hot-potato task passes one juggling ball between the hands, followed
by the preferred two-ball juggling task. One limb's action changes the next
limb's problem. Independently observed release and catch events, trajectories
and timing would test that coordination. Preserve uncertainty when an event or
body part is not observed, and compare with cue- and history-based controls.

\paragraph{Punch discrimination and NeuroBoxer.}
Independently labelled punch classes would test discrimination and movement
prediction before becoming the NeuroBoxer virtual game. Boxing includes hand
and wrist stabilization as well as arm motion. A prime-era Mike Tyson opponent
is a proposed virtual game character. A virtual opponent supplies no physical
opponent or evidence of safety. Rights-qualified expert footage can inform
movement and tactics, but supplies no expert neural, EMG or force recordings.

\paragraph{Toy archery and NeuroArcher.}
Toy archery would add explicit grasp, asymmetric two-hand roles, draw under
opposing loads, release and return. Its proposed arrows have small, cup-shaped
suction tips and are carried in an over-arm pouch with multiple arrows. NeuroArcher would use that coordination to act
on a virtual target. Success still requires a declared outcome and an
independent comparison, rather than agreement between simulated components.

The whole snap is a further coordinated showcase, from arm-out/palm-down rest
through elbow and forearm swing and turn, wrist and hand upward during the snap,
then downward return. The cup reach remains a simpler precision example. The
motion, contact and anatomy needed for either task are specific; a human-looking
animation cannot demonstrate a monkey's ability to perform the snap.

\subsection{Separate motion control, neural contribution and physical writing}

The first game milestone uses body motion to control a virtual body. A later
EEG-only milestone would ask whether permitted brain observations contribute to
the specified action. Withheld EMG and measured movement cannot become predictor
inputs through alignment, filtering or cue leakage. Eye activity, covert muscle
activity, residual movement and shared task information need their own controls.
Keep the original brain-only, body-only and paired comparisons, target-persistence
baseline, partition rules and success criteria.

Aim to make the declared intent-to-action mapping as close to one-to-one as
possible by measuring its intermediary. State the participant, session, condition,
allowed context, tolerated error and independent reference. Reduced ambiguity
would be a task-specific result; it would not identify a unique hidden intention
or establish the best achievable mapping. Reverse reconstruction of recorded
EEG features remains conditional and cannot deliver sensation.

The proposed later physical writer is a separate experiment. EEG and surface
EMG record, whereas neuromuscular electrical stimulation (NMES) and functional electrical stimulation (FES) denote stimulation approaches that can produce
contractions or functional movement. Sensation delivery needs its own writer or
sensory transducer. Device, clinical, consent and physical admission remain
separate, and physical spinal scope remains read-only. These writers and
actuators are unqualified and outside the existing acquisition/virtual-work
budget scope; no recording or intervention is reported here.

\subsection{Ask what remains after the help ends}

A proposed blindfolded condition changes visual feedback. It cannot independently
establish neural contribution, causal control or retained skill. Those claims
need matched comparisons, including a trainer-off test that preserves any movement
bridge needed to perform the task. Measure immediate performance, delayed
retention and transfer to specified unfamiliar conditions separately. Report
remaining support, uncertainty and alternative explanations.

The model's training and the trainee's learning also need separate evidence.
Qualified paired data could fit predictions; optional control-policy reinforcement
learning needs states, actions, transitions and rewards. A simulated reward is
not evidence of biological teaching. Final held-out outcomes remain outside
fitting, fine-tuning, rewards and model selection; frozen inference is distinct
from a separately reviewed adaptive update. The retained LINK comparison failed
both model-admission gates and all five benefit claims. These prospective tasks
do not revise that result or establish retained learning, tissue repair or a cure.

% Prospective apparatus and protocol; scientific and publication acceptance remain open.
\section{Build a brain--body system we can test}

The proposed architecture becomes a first experiment only when its measurements
are specified. Begin with one hand and one movement that can be measured well.
Record the brain and muscle signals, capture the movement, and compare the
virtual hand's prediction from one declared input side with that independent
measurement. The error identifies what the next experiment must resolve.

The apparatus would use non-invasive EEG, surface EMG, a calibrated
motion-capture relay-composite, an open musculoskeletal-simulation physics-engine
with a lifelike view, and an open multimodal model trained on eligible neural
and muscular structure and activity data. The EEG/EMG observation pair and the
movement/physics pair would meet inside one conditional learner. Intent-related
brain signals and muscle activity would supply observations; measured movement
would check what the body actually did. Predictions from either end must answer
to observations outside the model. Form constrains what could happen; measured
function tests the prediction.

\subsection{One movement. Two directions.}

The forward test uses EEG and declared task context to predict supported muscle
activity and movement on withheld recordings. The reverse test uses muscle
activity and movement to predict recorded EEG features within that task. Compare
both predictions with independent measurements. Agreement between the models
cannot substitute for that test.

Movement measures what the body did. It narrows the explanations of brain
signals that remain plausible without uniquely identifying intent or subjective
experience. Retain stimulus, task, behaviour and separately reported context.
Show uncertainty instead of displaying a generated neural pattern as the
observed mind.

Surface EMG records muscle electrical activity. A claimed incoming sensory or
feedback pathway needs evidence of its own. The virtual body cannot supply that
missing observation. Keep recordings, predictions and simulation identifiable
throughout the test.

\subsection{A bridge must account for the spine}

The spinal cord helps shape the response. Muscle output depends on sensory
input, movement context, local circuitry and descending input; some spinal
pathways can retain learned changes. Reproducing a specified spinal function
would therefore require more than forwarding its signal: a brain--body bridge
must account for the processing it aims to reproduce.

A changed response can reflect immediate reflex modulation, lasting spinal
plasticity or peripheral adaptation. Those explanations need separate tests.
A change during the task does not establish retained memory, and a retained
change does not by itself locate that memory. Human reflex conditioning and
constrained mouse adaptation supply evidence for different claims.~\cite{thompson2009retained,thompson2013retained,lavaud2024}
The human studies do not establish brain-independent learning; the mouse
preparation does not demonstrate unrestricted walking or reaching. Our human
comparisons use abstracts and selected results, discussion or methods passages.
Relevant mouse-study passages were also examined in the author manuscript;
equivalence to the final published edition and the complete methods and
statistics have not been established here.

The periphery adds a different question: did the sensory signal itself change?
Passive proprioceptive training supplies a task-specific peripheral comparison;
its psychophysical and spindle-afferent experiments were separate.~\cite{ackerley2022retained}
Our passive-training comparison uses abstracts and selected results, discussion
or methods passages. None of these studies demonstrates universal skill
transfer or complex memory stored in a peripheral axon.

The next comparison must choose a function we can measure. Does the controller
respond appropriately to the supported context? Does an adaptation persist
through the declared test? Could peripheral or mechanical changes explain the
outcome? Keep recordings, estimates and simulated state in distinct visible
layers. Surface EMG cannot supply a missing sensory-nerve or spinal channel;
an explanatory cutaway cannot register recorded neural sites.

Externalization remains the research aim: can a specified spinal function be
reproduced outside its usual pathway? Replacing a computation, stimulating a
pathway, controlling a simulated limb and restoring biological function would
answer different questions. Each needs its own task- and site-specific evidence;
physical spinal scope remains read-only. Anatomical access does not establish
neural function, increased excitability does not establish useful plasticity,
and no fixed percentage of a complex skill can be allocated to brain, spine
or periphery from these proposals.

\subsection{Reading a signal and writing an output}

The first prototype would test the recording and prediction route with EEG,
EMG and independent movement measurements. To deliver feedback or act on a
body, it would also need the separate sensory or action route specified in
Section~\ref{b2-output-interfaces}. That method distinguishes indirect sensory
feedback from direct neural writing and virtual control from robotic or later
biological actuation, with the device, measurement, safety and clinical evidence
each would require.

Physical sensory writers, transducers and body actuators remain excluded from
this acquisition and virtual-work plan, unqualified and unpriced. Their absence
defines what this first prototype could test. The consolidated interface method
also states the excluded procedures and permissions: the proposal authorizes
no stimulation, participant acquisition, biological actuation, new equipment
purchase or training.

\subsection{The first prototype has to earn the next one}

The lean tier tests whether the selected signals and calibrated movement
reference can be synchronized and used for a bounded task. The serious-prototype tier
improves acquisition, independent references and source-specific bidirectional
evaluation. Expand cohorts, tasks, instrumentation and compute in the research-lab
tier only when earlier evidence justifies the next step.

Every tier would include equipment, calibration, accessories, storage,
integration, appropriate safety and consent work, independent references and
contingency. Planning figures would be stated in USD, with current, dated vendor sources
and a dated basis for any currency conversion. Listed prices, engineering
allowances and items requiring quotes would remain distinct. Each budget
defines a proposed scope; it authorizes no purchase and promises no complete
neural reconstruction.

\subsection{Let the model earn the claim}

Which open recordings can support the next model comparison? Their rights,
version, coverage, clock, units and missing modalities determine the answer.
Training would proceed only under an independently admitted protocol, with test
participants, sessions and tasks kept outside model selection. Simple baselines
would give the added model a comparison to beat. Examining artifact leakage,
absent inputs, timing errors, physically invalid solutions and uncertainty would
keep that comparison answerable to its actual conditions. Earlier failures and
reserved recordings remain constraints on it. The new ambition must answer to
the evidence already kept.

Begin with constrained virtual hand or foot control. Robotics, remote operation,
unfamiliar effectors and extra limbs remain later research directions, each
requiring separate physical, latency and safety evidence. A successful movement
prediction supports a movement claim. Biological and perceptual claims require
the corresponding broader evidence.

Rights-cleared images and lifelike interactive models would make the apparatus and its predictions inspectable. The comparison should identify the observation, the prediction and the error, with the conditions that limit each claim. A convincing image cannot supply a missing measurement or demonstrate biological function. The next experiment must make those distinctions visible enough to test.

\section{Testing the Proposal}

A prediction becomes testable when the question names an observable outcome and an independent reference can check the implemented answer. The measured error then gives refinement a direction. Apply that reasoning to acquisition, architecture, training, body mechanics, visual models, interfaces and communication. Synthetic augmentation can support the work only with real ground-truth support; it cannot invent missing training evidence or replace independent held-out real measurements for biological claims. Preserve failed trials and the original thresholds so that a more attractive result cannot quietly redefine success.\par

Interoceptomics, the proposed study joining internal-body sensing with connection maps, may extend the programme to heartbeat, breathing and other internal conditions. It needs distinct measurements and claims; the movement-first rationale does not directly reveal all sensation or private experience. Useful limited interfaces remain worth investigating even if a universal translator never emerges.\par

Return to the guided hand. A correct tune would establish a performance under
specified help; a delayed, controlled trainer-off comparison would ask what the
trainee retained. Better measurements and body-grounded tests might connect
those two achievements. Current results establish neither arbitrary skill
teaching nor complete biological virtualization, and the current development
models fail their benefit gates. The next generation must account for what
moves, what help causes the improvement, what remains without that help and
where the explanation breaks. That is the experiment the ambition still owes
us.\par

\section{Definitions and Lexicon}

The \hyperref[programme-definitions]{Scientific Article definition and lexicon supplement} is derived from the shared Brain2Body registry identified by version and fingerprint at the start of that supplement. The fingerprint identifies the exact definition source; it is not a scientific score. Preferred names, historical lookup labels, proposed scope and implemented software retain their distinct meanings. The Simple English Edition has its own supplement using the same registry. Necessary specialist terms are also explained where used in this article.\par

\section{References and Evidence Status}

Primary scientific sources are linked beside the claims they support. The accompanying bibliography preserves verified earlier identities, including the organoid paper's correction, and adds the ground-truth case-study sources used here. A published experiment supplies different evidence from a dated developer disclosure or an observation of pinned source code; those distinctions govern the claims throughout this article.\par

The accompanying \href{research-evidence.json}{computational results and limits}
describe unpublished numerical and software studies. The companion
NeuroTranslation manuscript is also an unpublished working draft; neither
source is peer-reviewed experimental support. The historical conceptual figure
explains a relationship rather than supplying new biological evidence. The
edition note records the remaining scientific and publication limits.\par

